\documentclass{article}
\usepackage[margin=1in]{geometry} % margins
\usepackage[justification=centering]{caption}
%\usepackage{calrsfs}
\usepackage{amsmath}
\usepackage{amsfonts}
\usepackage{amssymb}
\usepackage{parskip}
\usepackage{graphicx}
\usepackage{enumerate}
\usepackage{enumitem}
\usepackage{amsthm}
\usepackage{float}
\usepackage{tikz-cd}
\usepackage{ytableau}

\renewcommand{\implies}{\Rightarrow}
\newcommand{\longimplies}{\Longrightarrow}
\newcommand{\N}{\mathbb{N}}
\newcommand{\Z}{\mathbb{Z}}
\newcommand{\Q}{\mathbb{Q}}
\newcommand{\R}{\mathbb{R}}
\newcommand{\C}{\mathbb{C}}
\newcommand{\D}{\mathbb{D}}
\renewcommand{\H}{\mathbb{H}}
\newcommand{\F}{\mathbb{F}}
\newcommand{\K}{\mathbb{K}}
\newcommand{\Lf}{\mathbb{L}}
\newcommand{\Sp}{\mathbb{S}}
\newcommand{\cl}[1]{\overline{#1}}
\renewcommand{\emptyset}{\varnothing}
\newcommand{\norm}[1]{\left\|#1\right\|}
\newcommand{\maxnorm}[1]{\left\|#1\right\|_{\text{max}}}
\newcommand{\opnorm}[1]{\left\|#1\right\|_{\mathrm{op}}}
\newcommand{\supnorm}[1]{\left\|#1\right\|_{\infty}}
\newcommand{\id}{\operatorname{id}}
\newcommand{\sspan}{\operatorname{span}}
\newcommand{\tr}{\operatorname{tr}}
\newcommand{\ch}{\operatorname{char}}
\newcommand{\poly}{\mathrm{P}}
\newcommand{\mtx}{\mathcal{M}}
\newcommand{\seq}{\mathbf{Seq}}
\newcommand{\vv}[1]{\mathbf{#1}}
\newcommand{\dirsum}{\oplus}
\newcommand{\rank}{\operatorname{rank}}
\newcommand{\im}{\operatorname{im}}

\newcommand{\cj}[1]{\overline{#1}}
\newcommand{\inprod}[2]{\left\langle#1,#2\right\rangle}
\renewcommand{\Re}{\operatorname{Re}}
\renewcommand{\Im}{\operatorname{Im}}

\newcommand{\res}{\operatorname{res}}

\newcommand{\floor}[1]{\left\lfloor #1 \right\rfloor}

% Theorem environments
\usepackage{amsthm}

\theoremstyle{definition} % The "plain" style italicizes all body text.
	\newtheorem{thm}{Theorem}
		\numberwithin{thm}{section} % Theorem numbers are determined by section.
	\newtheorem{lemma}[thm]{Lemma}
	\newtheorem{prop}[thm]{Proposition}
	\newtheorem{cor}[thm]{Corollary}
	\newtheorem{conj}[thm]{Conjecture}

\theoremstyle{definition} % The "definition" style does not italicize body text..
    \newtheorem{defn}[thm]{Definition}
	\newtheorem{example}[thm]{Example}
    \newtheorem{exercise}[thm]{Exercise}

\usepackage{mathtools}

\usetikzlibrary{decorations.markings,positioning,decorations.text}

\def\xr{5}
\def\yr{5}

\usepackage{titlesec}
\titleformat{\section}[hang]{\normalfont\fontsize{14}{0}\bfseries}{Problem \thesection\ }{0pt}{}

\begin{document}

\section{}

Let $(M, d)$ be a metric space and $C(M)$ the space of continuous functions
 $f: M \to \R$ with metric
\[
    D(f, g) = \sup_{x \in M}|f(x) - g(x)|.
\]
We say that a set of functions $\mathcal{F} \subset C(M)$ is pointwise 
equicontinuous if for every $x \in M$ and $\varepsilon > 0$ there is a 
$\delta$ (depending on $x$ and $\varepsilon > 0$) for which
\[
    |f(x) - f(y)| < \varepsilon
\]
for all $f \in \mathcal{F}$ and $y$ satisfying $d(x, y) < \delta$. 
If $M$ is compact show that a pointwise equicontinuous family of 
functions is a equicontinuous.

\begin{proof}
Fix $\varepsilon > 0$ and for each $y \in M$, take a $\delta_y > 0$ 
for which $|f(x) - f(y)| < \varepsilon/2$ for all $d(x, y) < \delta_y$ and
$f \in \mathcal{F}$.
The balls $B_y(\delta_y/2)$ for $y \in M$ form an open cover of the compact
$M$, so there is a finite subcover $\{B_{y_i}(\delta_{y_i}/2)\}$ for
$y_1, \dots, y_n \in M$. Let $\delta = 
\min\{\delta_{y_1}/2, \dots, \delta_{y_n}/2\} > 0$ and consider any 
$x, y \in M$ with $d(x, y) < \delta$. $x \in B_{y_i}(\delta_{y_i}/2)$
for some $i$ because the balls $B_{y_i}(\delta_{y_i}/2)$ cover $M$, and
we note that
\[
    d(y, y_i) \leq d(x, y) + d(x, y_i) < \delta_{y_i}
\] 
so $y \in B_{y_i}(\delta_{y_i})$ as well. We then have for all 
$f \in \mathcal{F}$,
\[
    |f(x) - f(y)| \leq |f(x) - f(y_i)| + |f(y) - f(y_i)| < \varepsilon
\]
because $d(x, y_i) < \delta_{y_i}$ and $d(y, y_i) < \delta_{y_i}$. Hence
$\mathcal{F}$ is equicontinuous.
\end{proof}

\section{}

Let $\{(X_i, di)\}_{i = 1}^\infty$ be a sequence of metric spaces. 
Let $X = \prod_{i = 1}^\infty X_i$ and for any $x = (x_i)_i, y = (y_i)_i
\in X$ define,
\[
    d(x, y) := \sum_{j = 1}^\infty 
    \frac{1}{2^j}\frac{d_j(x_j, y_j)}{1 + d_j(x_j, y_j)}
\]
\begin{enumerate}[label=(\alph*)]
    \item Show that $(X, d)$ is a metric space.
    \item Let $x^{(n)} \in X$ and $x \in X$. Show that $x^{(n)} \to x$ 
    in $(X, d)$ iff $x^{(n)}_i \to x_i$ in $(X_i, d_i)$ for all $i$.
    \item Consider the case $X_i = \R$ and $d_i(x, y) := |x - y|$ 
    for all $i$. Then $\ell^p \subset X$ for all $1 \leq p \leq \infty$. 
    Compare convergence in $\ell^p$ with convergence in $(X, d)$.
\end{enumerate}
\begin{proof}
\textit{(a)}. Firstly, we note that the general term of the series 
$d(x, y)$ is positive, so its convergence follows from the series being
bounded.
\[
    \sum_{j = 1}^\infty \frac{1}{2^j}\frac{d_j(x_j, y_j)}{1 + d_j(x_j, y_j)}
    \leq \sum_{j = 1}^\infty \frac{1}{2^j} = 1
\]
hence $d(x, y)$ is well-defined. We also see that $d(x, y) \geq 0$ with
equality if and only if each general term of the series is 0. That in turn
is equivalent to $d_j(x_j, y_j) = 0$ for all $j$, and that is equivalent to
$x = y$. $d(x, y) = d(y, x)$ follows by symmetry of each $d_j$.

For the triangle inequlity, we note that the function 
$f:[0, \infty) \to [0, \infty), x \mapsto x/(1 + x)$ has derivative 
$f'(x) = (1 + x)^{-2} > 0$ for all $x$ in its domain.
This means that $f$ is increasing, so for all 
$x, y, z \in X, j$,
\[
    \frac{d_j(x_j, y_j)}{1 + d_j(x_j, y_j)} = f(d_j(x_j, y_j)) \leq 
    f(d_j(x_j, z_j)) + f(d_j(y_j, z_j)) = 
    \frac{d_j(x_j, z_j)}{1 + d_j(x_j, z_j)} + 
    \frac{d_j(y_j, z_j)}{1 + d_j(y_j, z_j)}
\]
because $d_j(x_j, y_j) \leq d_j(x_j, z_j) + d_j(y_j, z_j)$. It then follows
that 
\[
    d(x, y) = \sum_{j = 1}^\infty \frac{1}{2^j}
    \frac{d_j(x_j, y_j)}{1 + d_j(x_j, y_j)} \leq 
    \sum_{j = 1}^\infty \frac{1}{2^j}
    \left(\frac{d_j(x_j, z_j)}{1 + d_j(x_j, z_j)} 
    + \frac{d_j(y_j, z_j)}{1 + d_j(y_j, z_j)}\right)
    = d(x, z) + d(y, z)
\]
so $d$ is a metric.

\textit{(b)}. $(\longimplies)$. 
Suppose $x^{(n)} \to x$ in $(X, d)$ 
and fix $i, \varepsilon > 0$. We can take $N > 0$ such that 
$d(x^{(n)}, x) < \varepsilon / 2^{i + 1}$ and $d(x^{(n)}, x) < 1$ 
for $n \geq N$, and for all such $n$, we get 
\[
    d_i(x^{(n)}_i, x_i) \leq \frac{2d_i(x^{(n)}_i, x_i)}{1 + d_i(x^{(n)}_i, x_i)}
    \leq 2^{i + 1}\sum_{j = 1}^\infty \frac{1}{2^j}
    \frac{d_j(x_j, y_j)}{1 + d_j(x_j, y_j)} 
    < 2^{i + 1} \cdot \varepsilon / 2^{i + 1} 
    = \varepsilon
\]
Thus $x_i^{(n)} \to x_i$ in $(X_i, d_i)$.

$(\Longleftarrow)$. 
Suppose $x^{(n)}_i \to x_i$ in $(X_i, d_i)$ for each $i$. 
Fix $\varepsilon > 0$ and take $k > 0$ for which $1/2^k < \varepsilon / 2$.
For each $1 \leq i \leq k$ we can take $N_i > 0$ such that 
$d_i(x^{(n)}_i, x_i) < \varepsilon / 2$ for all $n \geq N_i$, and we let 
$N = \max\{N_1, \dots, N_k\}$. Then for all $n \geq N$, we get 
\begin{align*}
    d(x^{(n)}, x)
    &= \sum_{j = 1}^\infty \frac{1}{2^j}
    \frac{d_j(x_j^{(n)}, x_j)}{1 + d_j(x_j^{(n)}, x_j)} \\
    &= \sum_{j = 1}^k \frac{1}{2^j} 
    \frac{d_j(x_j^{(n)}, x_j)}{1 + d_j(x_j^{(n)}, x_j)}
    + \sum_{j = k + 1}^\infty \frac{1}{2^j}
    \frac{d_j(x_j, y_j)}{1 + d_j(x_j, y_j)} \\
    &< \sum_{j = 1}^k \frac{1}{2^j} \frac{\varepsilon}{2}
    + \sum_{j = k + 1}^\infty \frac{1}{2^j} \\
    &< \frac{\varepsilon}{2}\sum_{j = 1}^\infty \frac{1}{2^j} 
    + \frac{\varepsilon}{2}\sum_{j = 1}^\infty \frac{1}{2^j} \\
    &= \varepsilon
\end{align*}
Thus $x^{(n)} \to x$ in $(X, d)$.

\textit{(c)}.
Note the following for all $x, y \in \ell^p$:
\[
    |x_i - y_i| \leq \left(\sum_{n = 1}^\infty |x_n - y_n|^p\right)^{1/p}
    = d_p(x, y)
\]
This implies the following:
\[
    d(x, y) = \sum_{j = 1}^\infty \frac{1}{2^j}
    \frac{|x_j - y_j|}{1 + |x_j - y_j|} \leq \sup_i|x_i - y_i| 
    \leq d_p(x, y)
\]
which shows that any sequence that converges in $\ell^p$ also converges
in $(X, d)$. The converse is not true however. The sequence 
$x^{(n)} = e_n = (0, \dots, 0, 1, 0, \dots)$ with 1 in the $n$-th entry
and zero elsewhere is convergent in $(X, d)$, since $x^{(n)}_i \to 0$
for each $i$. However, for any $n \neq m$, 
$d_p(x^{(m)}, x^{(n)}) = 2^{1/p}$ for $p \neq \infty$ 
and $1$ if $p = \infty$, so the sequence isn't even Cauchy.
\end{proof}

\section{}

Let $f: \R_\geq \to \R$ be a continuous function s.t. 
$\lim_{n\to\infty} f(nx) = 0$ for all $x \in \R$ (here
$\R_\geq := \{x \in \R : x \geq 0\}$). Prove that 
$\lim_{x \to \infty} f(x) = 0$.

Hint: apply the Baire Category theorem in the following steps:
\begin{enumerate}[label=(\alph*)]
    \item Let $\varepsilon > 0$ and define 
    $E_k := \{x : |f(nx)| \leq \varepsilon, \forall n \geq k\}$. 
    Show that $E_k$ is closed.
    \item Show that $\R_\geq = \bigcup_k E_k$.
    \item Conclude by the BCT that some $E_k$ 
    contains an interval $(a, b)$.
    \item Show that $\bigcup_{n \geq k}(na, nb) \leq (T, \infty)$ 
    for some $T > 0$.
    \item Conclude that $\limsup_{t \to \infty} |f(t)| \leq \varepsilon$ 
    and so $\lim_{t \to \infty} f(t) = 0$.
\end{enumerate}
\begin{proof}
\textit{(a).}
Let $\varepsilon > 0$ and define $E_k := \{x : |f(nx)| 
\leq \varepsilon, \forall n \geq k\}$. For any sequence $x_m \in E_k$
that converges to $x \in \R$, we have $nx_m \to nx$ as $m \to \infty$
for any $n \geq k$, and since $|f(nx_m)| \leq \varepsilon$
and by continuity, we have $\varepsilon \geq 
\lim_{m \to \infty} |f(nx_m)| = |f(nx)|$. Hence $x \in E_k$ so
$E_k$ is closed.

\textit{(b).}
Given $x \in \R_\geq$, we have $\lim_{n \to \infty} f(nx) = 0$, so
for some $k$, $|f(nx)| < \varepsilon$ whenever $n \geq k$.
Thus $x \in E_k$ and so $\R_\geq \subseteq \bigcup_k E_k$. The other
subset inclusion is obvious.

\textit{(c).}
By BCT, $\R_\geq$ being complete and a union of countably many closed
sets $E_k$ implies one of $E_k$'s has a non-empty interior. 
That is some $E_k$ contains an interval $(a, b)$.

\textit{(d).}
Note that for a large enough $n$, we have $n(b - a) > a$, which implies 
that for all $n$ large enough, $(na, nb) \cup ((n + 1)a, (n + 1)b)$ is
an interval. This shows that 
$(T, \infty) \subset \bigcup_{n \geq k} (na, nb)$ for some $T > 0$.

\textit{(e).}
For each $x \in (T, \infty)$, we have $x \in (na, nb)$ for some 
$n \geq k$, so $ny = x$ for some $y \in (a, b) \subseteq E_k$, 
hence $|f(x)| = |f(ny)| \leq \varepsilon$. 
Thus $f(x) \to 0$ as $x \to \infty$.
\end{proof}

\section{}

Let $S \subseteq \ell^\infty$ be the set of bounded sequences s.t.
\[
    \lim_{n \to \infty} |x_n| = 0
\]
if $(x_n)_n \in S$.
\begin{enumerate}[label=(\alph*)]
    \item Prove that $S$ is closed.
    \item Let $B$ be the unit ball in $\ell^\infty$. Show that 
    $S \cap B$ is a closed and bounded set but is not compact.
    \item Fix $x \in S$. Let $\mathcal{F}$ be the set of 
    $y \in \ell^\infty$ so that $|y_n| \leq |x_n|$ for all $n \geq 0$. 
    Show that $\mathcal{F}$ is compact.
\end{enumerate}
\begin{proof}
\textit{(a).}
Let $x^{(n)}$ be a sequence in $S$ that converges to $x \in \ell^\infty$.
Then $|x_k| \leq |x^{(n)}_k - x_k| + |x^{(n)}_k|$ for all $k, n$.
For any $\varepsilon > 0$, we take $N > 0, K > 0$ such that 
$\supnorm{x^{(N)} - x} < \varepsilon/2$ and $|x^{(N)}_k| < \varepsilon/2$
for all $k \geq K$. By our earlier estimate, $|x_k| < \varepsilon$,
hence $x_k \to 0$ as $k \to \infty$, and thus $x \in S$.

\textit{(b).}
$S \cap B$ is closed as an intersection of two closed sets and 
clearly bounded (contained in a unit ball). Notice that the sequences
$e_i = (0, \dots, 0, 1, 0, \dots)$, with 1 in the $i$-th entry and zeros
everywhere else, is in $S \cap B$. For any $i \neq j$, we have 
$\supnorm{e_i - e_j} = 1$, so no subsequence of the sequence $(e_i)_i$
is even Cauchy. Thus $S \cap B$ is not compact.

\textit{(c).}
Let $y^{(n)}$ be a sequence in $\mathcal{F}$
\end{proof}

\section{}

Let $C([0, 1])$ be the set of continuous functions $f: [0, 1] \to \R$ with
\[
    D(f, g) = \sup_{x \in [0, 1]}|f(x) - g(x)|  
\]
For any $f \in C([0, 1])$ and $0 < \alpha \leq 1$ define,
\[
    \norm{f}_\alpha = \sup_x|f(x)| + 
    \sup_{x, y}\frac{|f(x) - f(y)|}{|x - y|^\alpha}.
\]
Let $S := \{f : \norm{f}_\alpha \leq 1\}$. Show that $S$ is a compact set.
\begin{proof}
By Arzela-Ascoli theorem, it suffices to show that $S$ is closed, bounded,
and equicontinuous. Boundedness follows from $\supnorm{f} \leq 
\norm{f}_\alpha \leq 1$ for all $f \in S$. For equicontinuity, let
$\varepsilon > 0$ and $\delta = \sqrt[\alpha]{\varepsilon} > 0$.
For any $x, y \in [0, 1]$ with $|x - y| < \delta$ and $f \in S$, we have
\[
    |f(x) - f(y)| = |x - y|^\alpha 
    \sup_{x, y}\frac{|f(x) - f(y)|}{|x - y|^\alpha} 
    \leq |x - y|^\alpha < \delta^\alpha = \varepsilon
\]
so $S$ is equicontinuous. It remains to show that $S$ is closed.

Let $f_n \in S$ be a sequence that converges to $f$ (uniformly). For fixed
$x \neq y$
\end{proof}

\section{}

Let $X = [0, 1]^n$ be the unit cube in $\R^n$. 
Show that multivariate polynomials are dense in $C_b(X)$.
\begin{proof}
By the Stone-Weierstrass theorem, it suffices to show that 
the multivariate polynomials $\mathcal{P}$ in $X$ form a function
algebra that vanishes nowhere and separates points. Clearly, for any
polynomials $f, g \in \mathcal{P}$ and constant $c \in \R$, we have that
$f + g, cf, fg$ are also polynomials, so $\mathcal{P}$ is a 
function algebra. The constant function $1$ is a polynomial, so
$\mathcal{P}$ vanishes nowhere. For distinct points 
$(a_1, \dots, a_n) \neq (b_1, \dots, b_n) \in X$, take $i$ such that
$a_i \neq b_i$ and note that the polynomial $f(x_1, \dots, x_n) = 
x_i - a_i$ is zero at $(a_1, \dots, a_n)$ but nonzero at
$(b_1, \dots, b_n)$, so $\mathcal{P}$ separates points.
Hence $\mathcal{P}$ is dense.
\end{proof}

\section{}

Let $(b_{ij})^\infty_{i,j=1}$ be a doubly semi-infinite matrix so that
\[
    s := \sup_i \sum_{j = 1}^\infty |b_{ij}| < \infty  
\]
\begin{enumerate}[label=(\alph*)]
    \item For any $x \in \ell^\infty$ define a sequence 
    $F(x) = (F(x)_i)_i$ by
    \[
        F(x)_i = \sum_j b_{ij}x_j^2  
    \]
    Show that $F$ is a continuous mapping from $\ell^\infty$ 
    to $\ell^\infty$.
    \item Show that there is an $r > 0$ so that for any 
    $a \in \ell^\infty$ with $\supnorm{a} < r$ the equation
    \[
        x = a + F(x)  
    \]
    has a solution $x \in \ell^\infty$. Show that this solution 
    is unique among solutions that are
    in some sufficiently small ball in $\ell^\infty$.
\end{enumerate}
\begin{proof}
\textit{(a).}
The series that defines $F(x)_i$ converges absolutely:
\[
    \sum_j |b_{ij}x_j|^2 \leq \supnorm{x}^2 \sum_j |b_{ij}| < \infty
\]
so $F(x)$ is well-defined. Fix $a \in \ell^\infty$, let 
$\varepsilon > 0$, and take $M > 0$ such that $M > \norm{a}$. 
Define $\delta = \min\{1, \varepsilon/(s(2M + 1))\} > 0$,
let $x \in B_a(\delta)$, and note the bound for all $j$:
\[
    |x_j + a_j| \leq |x_i - a_j| + 2|a_j| 
    < \delta + 2\supnorm{a} \leq 2M + 1
\]
We then have for all $i$,
\begin{align*}
    |F(x)_i - F(a)_i|
    &= \left|\sum_j b_{ij}(x_j^2 - a_j^2)\right| \\
    &\leq \sum_j |b_{ij}| \cdot |x_j + a_j| \cdot |x_j - a_j| \\
    &\leq \left(\sup_i \sum_j |b_{ij}|\right)(2M + 1)\supnorm{x - a} \\
    &< s(2M + 1)\delta \\
    &\leq \varepsilon
\end{align*}
Thus $\supnorm{F(x) - F(a)} \leq \varepsilon$, and hence $F$ is continuous.

\textit{(b)}.
Let $\Phi_a: X \to X, x \mapsto a + F(x)$ for $a \in \ell^\infty$.
We find conditions on $a$ that make $F$ into a contraction. 
Let $a \in \ell^\infty$ with $\supnorm{a} < r/2$.
For any $x, y \in \ell^\infty$ with 
$\supnorm{x}, \supnorm{y} \leq r$, we have
\[
    |\Phi_a(x)_i - \Phi_a(y)_i| = |F(x)_i - F(y)_i| \leq 2sr\supnorm{x - y}  
\]
so any $0 < r < 1/(2s)$ works. Let $B := \cl{B}_0(r)$, and note 
that if $F$ can be restricted to a function $B \to B$, then we will
have that $F$ is a contraction by the above bound. To ensure of this,
we show that $F(x)$ maps $x \in B$ to $B$.
\[
    |\Phi_a(x)_i| = |a_i + F(x)_i| \leq 
    |a_i| + \sum_j |b_{ij}x_j^2| \leq r/2 + sr^2 < r
\]
so $\norm{\Phi_a(x)} \leq r$.

Now let $a \in \ell^\infty$ with $\supnorm{a} < r/2$. $\Phi_a$ has a unique
fixed point $x_0$ in $B$, which is then the unique solution to
\[
    x = \Phi_a(x) = a + F(x) 
\]
\end{proof}

\section{}

\section{}

Let $X$ be a Banach space and $F: X \to X$ a contraction. Define 
$G(x) := x + F(x)$. Show that $G$ is a bijection with continuous inverse.
\begin{proof}
Let $x_0$ be the unique fixed point of $F$, and $k < 1$ the constant
that makes $F$ a contraction.
Define $H_n: X \to X$ by $H_1(y) = y - F(x_0)$ and 
$H_{n + 1}(y) = y - F(H_n(y))$ for $n \geq 1$. We then let
$H: X \to X$ be $H(y) = \lim_{n \to \infty} H_n(y)$.
To prove that $H$ is well-defined it suffices to show that
$H_n(y)$ is a Cauchy sequence. Note the following bound
\[
    \norm{H_{n + 1}(y) - H_n(y)} = \norm{F(H_n(y)) - F(H_{n - 1}(y))} 
    \leq k\norm{H_n(y) - H_{n - 1}(y)} \leq \cdots \leq k^n\norm{y - 2x_0}
\]
Given $\varepsilon > 0$, take $N > 0$ such that
\[
    \frac{k^N}{1 - k}\norm{y - 2x_0} < \varepsilon
\]
For any $N \leq n \leq m$, we have
\begin{align*}
    \norm{H_m(y) - H_n(y)}
    &\leq \norm{H_m(y) - H_{m - 1}(y)} + 
    \cdots + \norm{H_{n + 1}(y) - H_n(y)} \\
    &\leq k^{m - 1}\norm{y - 2x_0} + \cdots + k^n\norm{y - 2x_0} \\
    &= k^n(1 + k + \cdots + k^{m - n - 1})\norm{y - 2x_0} \\
    &\leq k^N\sum_{i = 0}^\infty k^i \norm{y - 2x_0} \\
    &= \frac{k^N}{1 - k}\norm{y - 2x_0} < \varepsilon
\end{align*}
Hence $H$ is well-defined. For any $x \in X$, we have
\begin{align*}
    G \circ H(x) 
    &= G(\lim_{n \to \infty}H_n(x)) \\
    &= \lim_{n \to \infty} H_n(x) + F(H_n(x)) \\
    &= \lim_{n \to \infty} H_{n + 1}(x) + F(H_n(x)) \\
    &= \lim_{n \to \infty} x - F(H_n(x)) + F(H_n(x)) \\
    &= x
\end{align*}
thus $G \circ H = \id$. Also we have the following:
\[
    \norm{F(x) - F(H_n(x + F(x)))} \leq k\norm{x - H_n(x + F(x))}
\]
This allows us to show that $H \circ G = \id$
\begin{align*}
    \norm{H \circ G(x) - x}
    &= \lim_{n \to \infty} \norm{H_n(x + F(x)) - x} \\
    &= \lim_{n \to \infty} \norm{F(x) - F(H_{n - 1}(x + F(x)))} \\
    &= \lim_{n \to \infty} \norm{F(x) - F(H_n(x + F(x)))} \\
    &\leq \lim_{n \to \infty} k\norm{x - H_{n - 1}(x + F(x))} \\
    &= k\lim_{n \to \infty} \norm{H_n(x + F(x)) - x} \\
    &= k^m\lim_{n \to \infty} \norm{H_n(x + F(x)) - x} \\
    &\to 0 \text{ as } m \to \infty
\end{align*}
Thus $H \circ G(x) = x$. Therefore $H$ is the inverse of $G$.

Lastly we show that $H$ is continuous. Consider a sequence $y_m \in X$
that converges to $y$. For all $n, k$, we have 
\[
    \norm{H_n(y_m) - H_n(y)} \leq \norm{y_m - y} 
    + k\norm{H_{n - 1}(y_m) - H_{n - 1}(y)} = \cdots 
    = (1 + k + \cdots + k^{n - 1})\norm{y_m - y} 
    \leq \frac{\norm{y_m - y}}{1 - k}
\]
Taking $n \to \infty$ we see that 
$\norm{H(y_m) - H(y)} \leq \norm{y_m - y}/(1 - k) \to 0$ 
as $m \to \infty$. Hence $H(y_m) \to H(y)$, and so $H$ is continuous.
\end{proof}

\end{document}